%======================================================================
\section{Track FND --- Foundation and upstream}\label{sec:fnd}
%======================================================================

Every other track depends on a fork that builds reproducibly, ships, and keeps up with FreeCAD. Two
foundation decisions changed in the interview. BruceWorx now tracks \textbf{upstream main} instead of
waiting for the 26.3 tag. Fixes found while reviewing the UI now have an epic of their own.

\begin{epic}{FND-01}{Build \& CI on Linux and macOS arm64}{M0}{\stProg}{FCB-P0-CI-01..06, P0-01, P0-06..08, P0-10}
Pixi environments and lockfiles for both platforms, ccache in CI, and the self-hosted Apple-silicon runner.
\file{BUILDING-BRUCEWORX.md} and \file{.github/workflows/bruceworx-ci.yml} exist, and the Linux STEP check
and hosted macOS lock are merged (\#4).
\emph{Acceptance:} a fresh clone builds on both platforms using only the documented commands, and CI is green
on both.
\end{epic}

\begin{epic}{FND-02}{Packaging, branding \& releases}{M0}{\stProg}{FCB-P0-02, P0-03, P0-09, P0-11, P0-17, P0-18}
The version string already reads \code{1.1.4-bw.1}, and the unsigned arm64 \code{.dmg} script is merged
(\#3). Remaining: tag a release carrying the dmg plus a source tarball, and decide on signing and
notarisation before anyone outside the team installs it.
\textbf{The KEYSTONE splash ships} (your decision, 24 Oct): only the splash screen is branded
(\file{package/bruceworx/branding/}: \file{branding.xml}, \file{splash.html}, \file{render\_splash.sh} drawing the
version into the image); the development build's application and executable names stay FreeCAD's, so its
preferences, add-ons and paths are unchanged.
\textbf{Keystone.app} (your answers, 25 Oct: fully renamed, its own settings, built from the working tree):
\file{package/bruceworx/make-keystone.sh} builds \file{build/release} and installs it into a bundle. A
\file{branding.xml} of its own renames the window, menu bar, Dock and About box (\code{Application},
\code{WindowTitle}, \code{ExeName = Keystone}). \code{ExeName} also moves its settings to
\file{\textasciitilde/Library/Preferences/Keystone} and \file{Application Support/Keystone}, so it runs beside the
development build. The bundle is self-contained: \file{keystone\_relocate.py} copies the 534 libraries it needs from the
pixi environment and rewrites and re-signs 1\,453 files so that nothing points into \file{.pixi} or the build tree.
The result is 1.7\,GB in \file{/Applications}, version \code{1.1.4-bw.1+<commit>[-dirty]}, with an icon drawn from the
splash's arch. Checked by launching it from \file{/Applications}: it reports itself as Keystone, uses its own settings
folders, its own Python and the native Qt platform, and loads all 20 workbenches and BruceWorx.
\emph{Acceptance:} the tagged release installs on a clean Mac and opens a Fusion STEP export
(\code{LE-01001}, 200$\times$200$\times$8 envelope).
\end{epic}

\begin{epic}{FND-03}{Track upstream FreeCAD main}{M0}{\stTodo}{FCB-P1-16..21, P0-12 (replaced)}
\file{UPSTREAM-SYNC.md} currently says ``stay on 1.1.4, rebase onto 26.3 when it ships''. That changes to
continuous tracking:
\begin{itemize}
  \item \code{upstream-main} stays a fast-forward mirror. A nightly job rebases the \code{[BW]} series onto it
        and runs the BruceWorx tests plus the headless acceptance corpus (\eref{FND-05}).
  \item A conflict or red build opens an issue automatically. The \emph{conflict budget} is the number of
        upstream files our series touches. It is reported each night and must not grow without a reason
        recorded in the PR.
  \item Product releases are still cut from a tested point (\code{main} after a green nightly), not from
        arbitrary upstream commits.
\end{itemize}
\emph{Why:} the Assembly, TechDraw and FEM work that parity depends on lands upstream continuously.
Rebasing once a year makes every conflict expensive at once.
\emph{Acceptance:} the nightly rebase runs for two weeks; each conflict is resolved within a working day.
\end{epic}

\begin{epic}{FND-04}{Licences \& notices}{M0}{\stTodo}{FCB-P0-13}
FreeCAD is LGPL-2.1-or-later. The roadmap bundles more components: SheetMetal (LGPL-2.1-or-later), OCP, CadQuery and build123d
(Apache-2.0), possibly an OpenSCAD binary (GPL-2, shipped as a separate executable) and a fastener library
(licence to check \verify). Each needs a line in \file{NOTICE} and a third-party licence file in the app
bundle.
\emph{Acceptance:} every bundled component is listed; each release has a source tarball attached.
\end{epic}

\begin{epic}{FND-05}{Acceptance corpus \& test harness}{M0}{\stTodo}{FCB-P1-09, P3-08}
One corpus serves every epic: native and STEP versions of LE parts and assemblies, the 2-DOF arm, the quad
frame, and the benchmark models in this document. A single command runs every headless check.
Screenshots come from the existing before/after script. The renders in this document already come from
headless FreeCAD (\file{scripts/}), and the harness generalises that.
\emph{Acceptance:} \code{pixi run acceptance} runs all checks in CI in under 15 minutes.
\end{epic}

\begin{epic}{FND-06}{Fixes found in review}{M0}{\stDone}{(new)}
\begin{enumerate}
  \item \textbf{Linear Pattern direction 2 drops a row} (\S\ref{sec:bug}). \code{Spacings2} is declared
        \code{(\{\})}, which stores \code{[0.0]}. Fix the declaration, migrate saved files, add a test.
  \item \textbf{Workbench trim not applied.} The pack lists FEM, CAM, BIM, Robot and others as disabled
        (\file{bruceworx\_pack.py:118}), but the screenshot shows 13 tabs. The uncommitted
        \code{param\_fallback} change suggests the pack crashed on \code{int(None)} \verify.
  \item \textbf{Overlay model tree missing} on the screenshot, although it is on by default; root cause unknown \verify.
\end{enumerate}
\emph{Acceptance:} a $4\times3$ pattern gives 12 instances; only intended workbenches are visible; the tree shows on first run.
\textbf{Done 25 Oct.} Checked in the GUI on a fresh profile: the $4\times3$ pattern gives 12 solids; seven
workbenches are enabled and the workspace bar's \emph{More} menu lists only Assembly, Draft, Part, Sketcher and
Spreadsheet (no stock selector); the Browser shows as soon as a document is open (the empty Start page has it
hidden); the Report view is clean. Item~2 was the \code{int(None)} crash, fixed in \code{9991ddf}; item~3 was the
auto-hide default, now off.
\end{epic}

\begin{epic}{FND-07}{On by default, each with an off switch}{M0}{\stDone}{FCB-P1-01, P1-10}
Merged in \#9. It remains a rule for every epic below: a feature ships enabled and can be switched off on
Preferences~$\rightarrow$~BruceWorx, and switching it off restores stock behaviour.
\end{epic}

\begin{epic}{FND-08}{Gaps found in real use}{M1}{\stTodo}{(new, 24 Oct)}
Found on 24 Oct while the Leonardo actuator was modelled in your vault \emph{Vector} through the bridge
(\S\ref{sec:pdm07-plan}) and by the morning's GUI test agents. Each is fixed in the epic that owns it:
\begin{enumerate}[nosep]
  \item \textbf{No ``New part in vault'' command.} New items enter the vault only through \emph{Import Files into
        Vault}, until \eref{PDM-09}'s New Project wizard; a single numbered part should be one command
        (\eref{PDM-04}). \textbf{Fixed 25 Oct:} BruceWorx $\rightarrow$ Vault $\rightarrow$ \emph{New Part in Vault\ldots}
        (also in the Files tab's menu) asks for part or assembly, program and description and shows the number; the
        file (an App::Part holding a Body, with the vault's part properties) is numbered, checked in as WIP~01,
        checked out to you and opened with the Body active. The number is taken only once the file exists, so a
        failure burns none (\code{vault\_setup.new\_item}; GUI test \code{TestNewPartInVault}).
  \item \textbf{New vaults still write \code{-REV}} in the part-number format (the import proposed
        \code{LE-01-01001-001-A}), although \eref{PDM-03} moved revisions out of the P/N: the Create Vault engine
        still writes the old format. \textbf{Fixed 25 Oct:} Create Vault and Vault Settings write the P/N form
        (\code{LE-01-01001-001}); the wizard's ``first revision'' and \code{vault init --first-revision} are gone.
  \item \textbf{Create Assembly fails silently} when a spreadsheet tab is in front: it needs a 3D view. It should
        switch to the document's 3D view or say why (\eref{ASM-01}). \textbf{Fixed 25 Oct:} the assembly was made
        and set to edit, but marked active in the \emph{sheet's} view, so the 3D view had no active assembly and every
        assembly command stayed greyed. Create and Activate Assembly now bring the 3D view to the front first
        (making one if there is none; \code{[BW]} in \file{CommandCreateAssembly.py}; GUI test).
  \item \textbf{A face sketch's ``centre'' lands off-centre} (\eref{MOD-12}). \textbf{Checked 25 Oct, open:} a sketch
        made on a face (\emph{New Sketch} with the face selected) has its origin at the body origin projected onto the
        face, not at the face's centre: on a plate from (10,10) to (70,50) the sketch's (0,0) is at (0,0,8), while the
        face centre is (40,30,8). Fusion and SolidWorks do the same, so I left it; question for you below.
  \item \textbf{A counterbored Hole's preview points out of the plate} (\eref{MOD-05}). \textbf{Checked 25 Oct, not
        reproduced:} a counterbored Hole on a top-face and on a bottom-face sketch cuts into the plate with its head at
        that face, and the panel's preview overlay sits on the face (GUI test \code{TestBruceWorxGaps}). The case seen
        on 24 Oct was made through the bridge; if it comes back, the file it happened in is what is needed.
  \item \textbf{Parts reading the Parameters sheet of another file do not rebuild} when the sheet changes; each has
        to be touched and recomputed by hand. An \emph{Update from Parameters} command, and a stale-part marker,
        belong to \eref{MOD-02}.
  \item \textbf{An assembly in the file that holds the sheet its parts read} makes a document cycle (``The graph
        must be a DAG''). The gearbox now keeps the sheet in a root file (LE-02-01001) and the assembly in its own
        (LE-02-01002); Create Assembly should warn before it builds the cycle (\eref{ASM-01}).
  \item \textbf{A vault item's description cannot be changed} after it is created: the root file still reads
        ``NEMA 17 planetary gearbox 5:1'' (\eref{PDM-04}).
  \item \textbf{A failed import burns part numbers}: the steel parts' failure left LE-01-01004 to 01006 unused
        (\eref{PDM-04}).
  \item \textbf{An unused WIP item cannot be cancelled}: only released items retire, so the M3\texttimes8 screw
        (LE-01-01015), replaced by the M3\texttimes6 (LE-01-01016) before any release, stays a live WIP item. A
        \emph{Cancel item} for never-released items, journaled like a retire, belongs to \eref{PDM-04}.
  \item \textbf{No joint driver.} Gear and belt joints couple only inside a drag (the solver takes their zero
        when a solve starts), and a drag step follows the pull only part way: turning the sun one full turn took
        4\,320 steps and still ended at 353\textdegree{} (the ratio came out 5.0000, the planets' spin matched
        theory). A revolute or gear joint driven to an angle, from Python and in a motion study, belongs to
        \eref{ASM-10}; an interference check over the parts (24 here, 276 pairs, by hand today) to \eref{ASM-06}.
  \item \textbf{Opening parts to link them offered ten check-outs}: the assembly command's transaction reached each
        part as the vault wrote its state. Fixed on 24 Oct: only the active document offers one.
  \item \textbf{Fixed on 24 Oct:} an import now rewrites links to earlier imports and to work copies, and a
        material without a density no longer stops it.
\end{enumerate}
\emph{Acceptance:} \emph{New part in vault} makes a numbered WIP~01 part without an import; a new vault's format has
no \code{REV}; Create Assembly works, or says why, with a spreadsheet in front; a face sketch's centre and a
counterbored hole land where they should; each is covered by a GUI test.
\end{epic}
